\chapter{Randomised Algorithms}\label{ch:randomised}

\begin{center}\itshape
A coin cannot make a wrong algorithm right, but it can leave the worst case with no one to aim at:\\
the probability is over our coins, never over the input.\\
\upshape(the principle of randomised algorithms, in our words)
\end{center}

In March 1947 John von Neumann sent a handwritten letter to Robert Richtmyer, head of the Theoretical Division at Los Alamos. It outlined how to follow the neutrons in a
bomb not by solving the equations of their diffusion, which nobody could do, but by simulating many of them at random on the new electronic computer, the ENIAC, and counting what
happened. The idea had come out of conversations with Stanislaw Ulam, a mathematician who, as his colleague Nicholas Metropolis remembered, relaxed playing solitaire and
was stimulated by playing poker. It was Metropolis who suggested the name: Ulam had an uncle who would borrow money from relatives because he ``just had to go to Monte Carlo''.
The \emph{Monte Carlo method} was published under that name by Metropolis and Ulam in 1949.

Thirty years later the coin had moved from simulation into algorithms. In 1976 Michael Rabin published a short paper, \emph{Probabilistic algorithms}, with a primality test that
was wrong with tiny probability and much faster than anything known (Chapter~10), and Robert Solovay and Volker Strassen published another in 1977. Here the randomness is not
in the problem but in the algorithm, which flips coins of its own; the input may be chosen by an adversary, and the guarantee holds for every input, with probability taken over the coins only.

This chapter collects the most beautiful examples. Karger's minimum cut algorithm contracts random edges and is right surprisingly often. Polynomial identities can be tested by evaluating at random points,
and this tests whether a graph has a perfect matching. Throwing balls into bins explains hash tables, and asking two bins instead of one is exponentially better. And data streams too large to store can be
summarised in a few counters, if the counters are random. We use Tome~3 for probability: linearity of expectation (Theorem~2.5), Markov's and Chebyshev's inequalities (Proposition~4.1), and Chernoff and
Hoeffding bounds (Theorems~4.6 and~4.9). Randomised quicksort and selection were already in Chapter~\ref{ch:sorting}.

\begin{think}
You are given two programs that are supposed to compute the same polynomial of degree $100$ in $10$ variables, and you can only run them. Expanding is out of the question: there could be
$\binom{110}{10}\approx4.6\times10^{13}$ coefficients. How would you test them, and how confident could you be?
\end{think}

\section{Las Vegas and Monte Carlo}\label{sec:lvmc}

A randomised algorithm reads, in addition to its input, a sequence of random bits. A \emph{Las Vegas} algorithm always answers correctly, and its running time is a random variable, whose
expectation we bound for every input; randomised quicksort is one. A \emph{Monte Carlo} algorithm always runs within a fixed time bound, and may answer wrongly with a probability that
we bound for every input. A decision algorithm has \emph{one-sided error} if one of its answers is always right, for example ``no'' when it says ``no''. (The name Las Vegas was proposed by L\'aszl\'o Babai in 1979 to
contrast with Monte Carlo.)

\begin{proposition}[amplification]\label{prop:amplify}
\begin{enumerate}
\item A Monte Carlo algorithm with one-sided error which errs with probability at most $1-p$, repeated $k$ times independently (answering ``yes'' if one run says ``yes'', when ``yes'' is the always-correct answer), errs with probability at most $(1-p)^k\le e^{-pk}$.
\item A Monte Carlo algorithm which answers correctly with probability at least $\tfrac12+\varepsilon$, repeated $k$ times independently with a majority vote, errs with probability at most $e^{-2\varepsilon^2k}$.
\item A Las Vegas algorithm with expected time $T$, stopped after $2T$ steps (answering ``don't know''), fails with probability at most $\tfrac12$.
\end{enumerate}
\end{proposition}
\begin{proof}
(1) The combined answer is wrong only if all $k$ independent runs are, and $1-p\le e^{-p}$. (2) Let $X_i=1$ if run $i$ is correct; $S=\sum X_i$ has $\mathbb ES\ge(\tfrac12+\varepsilon)k$, and the majority is wrong only if $S\le k/2$, that is
$S-\mathbb ES\le-\varepsilon k$, of probability at most $e^{-2\varepsilon^2k}$ by Hoeffding's inequality (Tome~3, Theorem~4.9). (3) By Markov's inequality, the running time exceeds $2T$ with probability at most $\tfrac12$.
\end{proof}

So the constants do not matter: an error probability of $\tfrac13$ becomes $10^{-30}$ after a few hundred repetitions, far below the probability of a hardware fault during the computation.

\section{Minimum cuts by contraction}\label{sec:karger}

A \emph{cut} of a connected multigraph $G=(V,E)$ is a partition of $V$ into two nonempty parts, and its \emph{size} is the number of edges between them; we want a cut of minimum size $k$. Chapter~\ref{ch:flows}
finds one with $n-1$ maximum flow computations (fix $s$, try every $t$). In 1993 David Karger, then a graduate student at Stanford, proposed something that looks too naive to work.

To \emph{contract} an edge $uv$ is to merge $u$ and $v$ into one vertex, keeping all the other edges (parallel edges included) and deleting the loops that the edges between $u$ and $v$ become. Karger's
\emph{contraction algorithm} repeats: pick an edge uniformly at random among the remaining ones and contract it, until two vertices remain; it outputs the cut formed by the two sets of original vertices
merged into them (Figure~\ref{fig:karger}).

\begin{figure}[ht]
\centering
\fig{ch08-karger}
\caption{A lucky run of the contraction algorithm. The minimum cut, of size $2$ (in red), separates $\{a,b,c\}$ from $\{d,e\}$; no red edge is ever contracted, and it is what remains at the end.}
\label{fig:karger}
\end{figure}

\begin{theorem}[Karger, 1993]\label{thm:karger}
For every minimum cut $C$ of a connected multigraph with $n\ge2$ vertices, the contraction algorithm outputs $C$ with probability at least $\binom n2^{-1}=\frac2{n(n-1)}$.
\end{theorem}
\begin{proof}
Every cut of a contracted multigraph is a cut of $G$ (a vertex of the contracted graph is a set of original vertices), so contracted graphs have minimum cut at least $k$, and every vertex has degree at least $k$: a multigraph with
$i$ vertices has at least $ik/2$ edges. As long as no edge of $C$ has been contracted, the two sides of $C$ are unions of current vertices, and all $k$ edges of $C$ are still present (loops only come from edges inside a merged set). So,
given that $C$ has survived until $i$ vertices remain, the next contraction destroys it with probability at most $k/(ik/2)=2/i$, and
\[
P(C\text{ survives})\ge\prod_{i=3}^n\Bigl(1-\frac2i\Bigr)=\prod_{i=3}^n\frac{i-2}i=\frac{(n-2)!\,2!}{n!}=\frac2{n(n-1)}.
\]
If $C$ survives until two vertices remain, these are its two sides, and the output is $C$.
\end{proof}

\begin{corollary}\label{cor:fewmincuts}
A connected multigraph with $n$ vertices has at most $\binom n2$ minimum cuts.
\end{corollary}
\begin{proof}
The events ``the output is $C$'', for the different minimum cuts $C$, are disjoint, and each has probability at least $\binom n2^{-1}$.
\end{proof}

The cycle $C_n$ has exactly $\binom n2$ minimum cuts (Exercise~\ref{ex:cutsalpha}), so the bound is sharp: a probabilistic proof of a purely combinatorial fact. Repeating the algorithm $\binom n2\ln n$ times and keeping the smallest cut found errs with probability at most $1/n$ by
Proposition~\ref{prop:amplify}; with an adjacency matrix of multiplicities, one run costs $O(n^2)$ (pick a vertex with probability proportional to its degree, then a neighbour proportional to the multiplicity, which picks a uniform edge in $O(n)$; merge two rows in $O(n)$), and the total is $O(n^4\log n)$, worse than flows. The fix, found by Karger and Clifford Stein in 1993, starts from an observation: the early contractions are safe, and only the last ones are risky. So
share the safe part and repeat only the risky part.

The \emph{recursive contraction algorithm} $\mathrm{RC}(G)$: if $G$ has $n\le6$ vertices, return a minimum cut by trying all $2^{n-1}-1$ cuts. Otherwise let $t=\lceil1+n/\sqrt2\,\rceil$, and twice, independently, contract $G$ at random down to $t$ vertices and call $\mathrm{RC}$ on the result;
return the smaller of the two cuts.

\begin{theorem}[Karger and Stein, 1993]\label{thm:kargerstein}
$\mathrm{RC}$ runs in time $O(n^2\log n)$ and returns a minimum cut with probability at least $1/(2\log_2n+8)$. Repeated $O(\log^2n)$ times, it finds a minimum cut with probability at least $1-1/n$, in time $O(n^2\log^3n)$.
\end{theorem}
\begin{proof}
\emph{Depth.} Let $c_0=2/(1-1/\sqrt2)\approx6.83$, so that $2-c_0=-c_0/\sqrt2$. For $n\ge7$, $t\le n/\sqrt2+2$, hence $t-c_0\le(n-c_0)/\sqrt2$: along a branch of the recursion, the sizes $n_j$ at depth $j$ satisfy
$n_j-c_0\le(n-c_0)2^{-j/2}$. The recursion stops when $n_j\le6$, which happens as soon as $(n-c_0)2^{-j/2}<0.17$. As $2\log_2(1/0.17)<5.2$, the depth $D(n)$ is at most $2\log_2n+7$ (and $D(n)=0$ for $n\le6$).

\emph{Time.} With an adjacency matrix, contracting from $n_j$ to $n_{j+1}$ vertices costs $O(n_j^2)$, and there are $2^j$ calls at depth $j$, each on at most $c_0+n2^{-j/2}$ vertices. The work at depth $j$ is
$O\bigl(2^j(c_0^2+n^22^{-j})\bigr)=O(n^2+2^j)$, and summing over $j\le D(n)$ gives $O(n^2D(n)+2^{D(n)})=O(n^2\log n)$, since $2^{D(n)}\le128n^2$. The base cases cost $O(1)$ each.

\emph{Success.} We show by induction on the depth $d=D(n)$ that $\mathrm{RC}$ returns a minimum cut with probability at least $1/(d+1)$. For $d=0$ it is exact. Otherwise fix a minimum cut $C$, of size $k$. As in the proof of
Theorem~\ref{thm:karger}, $C$ survives the contraction from $n$ to $t$ vertices with probability at least
\[
\prod_{i=t+1}^n\Bigl(1-\frac2i\Bigr)=\frac{t(t-1)}{n(n-1)}\ge\frac{(1+n/\sqrt2)(n/\sqrt2)}{n^2}\ge\frac12.
\]
If it survives, the contracted graph has minimum cut exactly $k$, and any minimum cut of it is a cut of size $k$ of $G$. Its depth is $d-1$, so by induction the recursive call then returns a minimum cut with probability at least
$q=1/d$. So each of the two independent branches succeeds with probability at least $q/2$, and
\[
P(\text{success})\ge1-\Bigl(1-\frac q2\Bigr)^2=q-\frac{q^2}4.
\]
The function $x\mapsto x-x^2/4$ is increasing on $[0,2]$, and
$\frac1d-\frac1{4d^2}\ge\frac1{d+1}$ because $\frac1d-\frac1{d+1}=\frac1{d(d+1)}\ge\frac1{4d^2}$. This completes the induction; with $d\le2\log_2n+7$ we get the stated bound. Finally, $r$ independent repetitions all fail with probability at most
$(1-\frac1{2\log_2n+8})^r\le e^{-r/(2\log_2n+8)}$, which is at most $1/n$ for $r=\lceil(2\log_2n+8)\ln n\rceil=O(\log^2n)$.
\end{proof}

\begin{think}
Why $n/\sqrt2$? Show that, with $t=n/\alpha$, the survival probability of a fixed minimum cut is about $1/\alpha^2$, and that the recursion keeps a success probability of order $1/\log n$ only if two branches compensate this loss, that is $\alpha\le\sqrt2$.
What happens to the running time if $\alpha<\sqrt2$?
\end{think}

\section{Testing identities}\label{sec:identities}

A polynomial that is not zero cannot vanish at many points: this simple fact is the heart of a whole family of randomised algorithms.

\begin{lemma}[Schwartz--Zippel; DeMillo and Lipton, 1978; Zippel, 1979; Schwartz, 1980]\label{lem:schwartzzippel}
Let $f\in K[x_1,\dots,x_v]$ be a nonzero polynomial of total degree at most $d$ over a field $K$, and $S\subseteq K$ a finite set. If $r_1,\dots,r_v$ are independent and uniform in $S$, then
$P\bigl(f(r_1,\dots,r_v)=0\bigr)\le d/|S|$.
\end{lemma}
\begin{proof}
By induction on $v$. For $v=1$, a nonzero polynomial of degree at most $d$ has at most $d$ roots in a field. For $v\ge2$, write $f=\sum_{i=0}^kx_1^ig_i(x_2,\dots,x_v)$ with $g_k\ne0$; then $g_k$ has total degree at most $d-k$, so by induction
$P(g_k(r_2,\dots,r_v)=0)\le(d-k)/|S|$. When $g_k(r_2,\dots,r_v)\ne0$, the polynomial $f(x_1,r_2,\dots,r_v)$ in $x_1$ is nonzero of degree $k$, and $r_1$ is independent of $r_2,\dots,r_v$, so it is a root with probability at most $k/|S|$. Altogether
$P(f(r)=0)\le\frac{d-k}{|S|}+\frac k{|S|}$.
\end{proof}

So to test whether two programs compute the same polynomial of degree $d$, evaluate both at a random point of $S^v$ with $|S|\ge2d$, and repeat: each test that agrees halves the probability of a false alarm. Nobody knows a
\emph{deterministic} polynomial algorithm for this problem, \emph{polynomial identity testing}, when the polynomials are given by arithmetic circuits; Kabanets and Impagliazzo showed in 2004 that finding one would prove lower bounds in circuit complexity that seem far out of reach (Tome~C3).

The most striking application concerns matchings. Let $G$ be a graph on the vertices $1,\dots,n$, and give each edge $e=ij$ ($i<j$) an indeterminate $x_e$. The \emph{Tutte matrix} $T$ of $G$ is the $n\times n$ matrix with
$T_{ij}=x_e$ and $T_{ji}=-x_e$ for each edge $e=ij$ with $i<j$, and $0$ elsewhere.

\begin{theorem}[Tutte, 1947]\label{thm:tuttematrix}
$G$ has a perfect matching if and only if $\det T$ is not the zero polynomial.
\end{theorem}
\begin{proof}
Expand $\det T=\sum_\sigma\operatorname{sgn}(\sigma)\prod_iT_{i\sigma(i)}$ over the permutations of $\{1,\dots,n\}$. A term is nonzero only if every $\{i,\sigma(i)\}$ is an edge; in particular $\sigma$ has no fixed point.

If $M$ is a perfect matching, consider the monomial $\prod_{e\in M}x_e^2$. A term $\prod_iT_{i\sigma(i)}$ equals $\pm$ this monomial only if every edge $\{i,\sigma(i)\}$ is in $M$, that is, $\sigma(i)$ is the partner of $i$ in $M$: $\sigma$ is the product of the
$n/2$ transpositions of $M$. Its term is $(-1)^{n/2}\prod_{e\in M}\bigl(x_e\cdot(-x_e)\bigr)=\prod_{e\in M}x_e^2$, so this monomial has coefficient $1$ in $\det T$, which is not zero.

Conversely, suppose $\det T\ne0$. Call a permutation \emph{odd-cyclic} if it has a cycle of odd length, and pair each odd-cyclic $\sigma$ with nonzero term with the permutation $\sigma'$ obtained by reversing the odd cycle that contains the
smallest vertex lying on an odd cycle. This cycle has length at least $3$, so $\sigma'\ne\sigma$; $\sigma'$ has the same cycles as sets, so $(\sigma')'=\sigma$ and $\operatorname{sgn}\sigma'=\operatorname{sgn}\sigma$; and since $T_{ji}=-T_{ij}$, reversing a cycle of odd length
multiplies its product by $(-1)^{\text{odd}}=-1$. So the terms of $\sigma$ and $\sigma'$ cancel, and the odd-cyclic permutations contribute nothing in total. As $\det T\ne0$, some $\sigma$ with nonzero term has only cycles of even length. Each such cycle
$i_1\to i_2\to\dots\to i_{2\ell}\to i_1$ runs along edges of $G$, and the edges $i_1i_2,i_3i_4,\dots,i_{2\ell-1}i_{2\ell}$ match its vertices (for $\ell=1$ it is a single edge). Together they form a perfect matching.
\end{proof}

\begin{corollary}[Lov\'asz, 1979]\label{cor:lovasz}
Let $p\ge2n$ be a prime. Substituting independent uniform elements of $\mathbb F_p$ for the $x_e$ and computing the determinant modulo $p$ decides whether $G$ has a perfect matching in time $O(n^3)$, never answering ``yes'' wrongly and
answering ``no'' wrongly with probability at most $\tfrac12$.
\end{corollary}
\begin{proof}
The determinant is a polynomial with integer coefficients of total degree $n$, and the proof of Theorem~\ref{thm:tuttematrix} shows that if $G$ has a perfect matching it has a coefficient equal to $1$, so it is nonzero modulo $p$. If there is no perfect matching, all the
terms cancel identically, also modulo $p$, and the evaluation is $0$. In the first case Lemma~\ref{lem:schwartzzippel} with $S=\mathbb F_p$ bounds the probability of a zero by $n/p\le\frac12$. A determinant modulo $p$ costs $O(n^3)$ field operations by Gaussian elimination.
\end{proof}

This is slower than Edmonds' algorithm of Chapter~\ref{ch:flows}, but it is \emph{algebra}: determinants can be computed in parallel in time $O(\log^2n)$ with polynomially many processors, so perfect matchings can be decided in randomised parallel polylogarithmic time. Finding a matching in parallel needs one more
idea, the \emph{isolation lemma} of Mulmuley, Vazirani and Vazirani (1987): random weights make the minimum weight perfect matching unique with good probability. Deterministic parallel algorithms are known only with quasi-polynomially many processors (Fenner, Gurjar and Thierauf in 2016 for bipartite graphs, Svensson and Tarnawski in 2017 in general); polylogarithmic time with polynomially many processors is open.
Exercise~\ref{ex:tuttefind} finds the matching sequentially.

\begin{think}
Is the hypothesis ``$p\ge2n$'' really needed? What can go wrong when $p$ is small, say $p=2$?
\end{think}

\emph{Fingerprints.} The same idea checks computations. In 1977 R\=usi\c n\v s Freivalds observed that a claimed matrix product can be verified much faster than it can be computed.

\begin{proposition}[Freivalds, 1977]\label{prop:freivalds}
Let $A,B,C$ be $n\times n$ matrices over a field. If $AB\ne C$ and $r$ is uniform in $\{0,1\}^n$, then $P(A(Br)=Cr)\le\tfrac12$. The test costs $O(n^2)$ operations.
\end{proposition}
\begin{proof}
Let $D=AB-C\ne0$, with $d_{ij}\ne0$. Fix the coordinates of $r$ other than $r_j$. Then $(Dr)_i=d_{ij}r_j+\gamma$ for a constant $\gamma$, which vanishes for at most one of the two values of $r_j$. So $P((Dr)_i=0)\le\tfrac12$, and
$Dr=0$ is even less likely. The vectors $Br$, $A(Br)$ and $Cr$ cost three matrix--vector products.
\end{proof}

\begin{proposition}[polynomial fingerprints]\label{prop:fingerprint}
Let $p$ be a prime, and for a string $a=a_0a_1\dots a_{n-1}$ of elements of $\mathbb F_p$ let $F_x(a)=\sum_ia_ix^i\in\mathbb F_p$. If $a\ne b$ have length $n$ and $x$ is uniform in $\mathbb F_p$, then $P(F_x(a)=F_x(b))\le(n-1)/p$.
\end{proposition}
\begin{proof}
$F_x(a)-F_x(b)$ is the value at $x$ of the nonzero polynomial $\sum_i(a_i-b_i)X^i$, of degree at most $n-1$.
\end{proof}

Two computers holding files of $n$ bytes can compare them by exchanging one fingerprint, $O(\log n)$ bits, with error at most $1/n$ if $p>n^2$ (Exercise~\ref{ex:stringeq}): any deterministic protocol needs $n$ bits for this. The fingerprints of all the windows of length $m$ of a text can be updated in
$O(1)$ each as the window slides, which is the Karp--Rabin pattern matching algorithm of Chapter~11.

\section{Balls into bins}\label{sec:bins}

Throw $n$ balls independently and uniformly into $n$ bins. This is the model of a hash table with chaining, of a load balancer sending jobs to random servers, of birthdays in a large class. The average load is $1$; how large is the largest?

\begin{theorem}\label{thm:maxload}
Let $n\ge3$ and $L$ be the maximum load. Then $P\bigl(L\ge3\ln n/\ln\ln n\bigr)\le1/n$ and $P\bigl(L<\lfloor\ln n/(2\ln\ln n)\rfloor\bigr)\le4/\sqrt n$.
\end{theorem}
\begin{proof}
Let $L_i$ be the load of bin $i$. \emph{Upper bound.} Let $k=\lceil3\ln n/\ln\ln n\rceil$. Bin $1$ receives at least $k$ balls only if some set of $k$ balls all land in it, so by the union bound and $k!\ge(k/e)^k$,
\[
P(L_1\ge k)\le\binom nk\frac1{n^k}\le\frac1{k!}\le\Bigl(\frac ek\Bigr)^k.
\]
Put $u=\ln\ln n>0$. Then $\ln(k/e)\ge\ln3+\ln\ln n-\ln u-1$, and the function $u\mapsto u/3-\ln u$ has its minimum $1-\ln3$ at $u=3$, so $\ln(k/e)\ge\tfrac23u$, and $k\ln(k/e)\ge\frac{3\ln n}u\cdot\frac{2u}3=2\ln n$. Hence $P(L_1\ge k)\le n^{-2}$,
and by the union bound over the $n$ bins, $P(L\ge k)\le1/n$.

\emph{Lower bound.} Let $k=\lfloor\ln n/(2\ln\ln n)\rfloor$; if $k=0$ there is nothing to prove. Let $X$ be the number of bins with load at least $k$. First, using $\binom nk\ge(n/k)^k$ and $(1-1/n)^n\ge\frac14$,
\[
P(L_1\ge k)\ge\binom nk\frac1{n^k}\Bigl(1-\frac1n\Bigr)^{n-k}\ge\frac1{4k^k},
\]
and $k\ln k\le\frac{\ln n}{2\ln\ln n}\ln\ln n=\frac12\ln n$ (as $k\le\ln n$), so $\mathbb EX=nP(L_1\ge k)\ge\sqrt n/4$.

Next, the indicators of the events $\{L_i\ge k\}$ are negatively correlated. Given $L_i=a$, the other $n-a$ balls are uniform among the $n-1$ other bins, so $L_j$ is binomial with parameters $n-a$ and $\frac1{n-1}$, and $g(a)=P(L_j\ge k\mid L_i=a)$ is
nonincreasing in $a$ (fewer balls can only lower it, as one sees by throwing the balls one at a time). So $P(L_j\ge k\mid L_i\ge k)\le g(k)\le P(L_j\ge k\mid L_i<k)$, and since $P(L_j\ge k)$ is a weighted average of these two conditional probabilities, it is
at least the first: $P(L_i\ge k,L_j\ge k)\le P(L_i\ge k)P(L_j\ge k)$. Then $\operatorname{Var}X$ is the sum of the variances of the indicators, each at most its mean, plus nonpositive covariances, so $\operatorname{Var}X\le\mathbb EX$. By Chebyshev's inequality,
$P(X=0)\le P(|X-\mathbb EX|\ge\mathbb EX)\le\operatorname{Var}X/(\mathbb EX)^2\le1/\mathbb EX\le4/\sqrt n$.
\end{proof}

So the maximum load is of order $\ln n/\ln\ln n$; a finer analysis shows that it is $(1+o(1))$ times $\ln n/\ln\ln n$ with high probability. For a million balls, about $8$ or $9$. Now change the rule slightly: each ball chooses \emph{two} bins at random and goes into the
less loaded one (ties broken in any way). Figure~\ref{fig:loads} shows the effect.

\begin{figure}[ht]
\centering
\fig{ch08-loads}
\caption{Number of bins with each load, for $100\,000$ balls in $100\,000$ bins (logarithmic scale). With one random choice the maximum load is $7$; with the less loaded of two choices it is $3$.}
\label{fig:loads}
\end{figure}

To analyse two choices we need a Chernoff bound for sums of dependent indicators, each of which is unlikely whatever happened before.

\begin{lemma}\label{lem:chernoffdom}
Let $Y_1,\dots,Y_n\in\{0,1\}$ be produced by a random process in steps, $Y_t$ being determined at step $t$, and suppose that, whatever happened in the first $t-1$ steps, the conditional probability that $Y_t=1$ is at most $p$. Let $S=\sum_tY_t$ and $\mu\ge np$. Then
$P(S\ge2\mu)\le e^{-\mu/3}$.
\end{lemma}
\begin{proof}
For $\lambda\ge0$, conditioning on the first $t-1$ steps, $\mathbb E[e^{\lambda Y_t}\mid\text{past}]=1+P(Y_t=1\mid\text{past})(e^\lambda-1)\le1+p(e^\lambda-1)\le e^{p(e^\lambda-1)}$. Multiplying by $e^{\lambda(Y_1+\dots+Y_{t-1})}$, which is determined by the past, and taking expectations,
$\mathbb E[e^{\lambda(Y_1+\dots+Y_t)}]\le e^{p(e^\lambda-1)}\mathbb E[e^{\lambda(Y_1+\dots+Y_{t-1})}]$, so $\mathbb E[e^{\lambda S}]\le e^{np(e^\lambda-1)}\le e^{\mu(e^\lambda-1)}$. Markov's inequality with $\lambda=\ln2$ gives $P(S\ge2\mu)\le e^{-2\mu\ln2}e^{\mu}=e^{-(2\ln2-1)\mu}\le e^{-\mu/3}$.
\end{proof}

This is the moment generating function argument of Tome~3, Theorem~4.6, which never needed independence, only the bound on each conditional probability (Lemma~\ref{lem:dominate} gave a coupling version of the same idea).

\begin{theorem}[Azar, Broder, Karlin and Upfal, 1994]\label{thm:twochoices}
Throw $n\ge3$ balls into $n$ bins, each ball choosing two bins independently and uniformly and going into the less loaded one. Then the maximum load is at most $\log_2\log_2n+6$ with probability at least $1-145(\ln n)^2/n$.
\end{theorem}
\begin{proof}
The \emph{height} of a ball is the load of its bin just after it is placed. Let $\nu_i(t)$ be the number of bins of load at least $i$, and $\mu_i(t)$ the number of balls of height at least $i$, after $t$ balls. A bin of load at least $i$ contains a ball of
height exactly $i$, so $\nu_i(t)\le\mu_i(t)$; and $\nu_i(t)$ grows with $t$. A ball gets height at least $i+1$ only if \emph{both} its choices have load at least $i$, which, given the past, happens with probability $(\nu_i(t-1)/n)^2$.

Define $\beta_4=n/4$ and $\beta_{i+1}=2\beta_i^2/n$, so that $2\beta_{i+1}/n=(2\beta_i/n)^2$ and $\beta_{4+j}=\frac n2\,2^{-2^j}$. Let $E_i$ be the event $\nu_i(n)\le\beta_i$. The event $E_4$ always holds, since $n$ balls fill at most $n/4$ bins up to load $4$. Let $i^*$ be the least
$i\ge4$ with $\beta_i^2/n<6\ln n$. For $j=\lceil\log_2\log_2n\rceil$, $2^{2^j}\ge n$ and $\beta_{4+j}\le\frac12$, so $i^*\le4+\lceil\log_2\log_2n\rceil$.

\emph{The induction.} Let $4\le i<i^*$ and $p_i=(\beta_i/n)^2$, so $np_i=\beta_i^2/n\ge6\ln n$. Let $Y_t=1$ if ball $t$ has height at least $i+1$ and $\nu_i(t-1)\le\beta_i$. Whatever the past, $P(Y_t=1)\le p_i$. On $E_i$, every $\nu_i(t-1)\le\beta_i$, so $\sum_tY_t=\mu_{i+1}(n)\ge\nu_{i+1}(n)$.
Hence, by Lemma~\ref{lem:chernoffdom} with $\mu=np_i$,
\[
P(E_i\cap\overline{E_{i+1}})\le P\Bigl(\sum_tY_t>2np_i\Bigr)\le e^{-np_i/3}\le n^{-2},
\]
since $\beta_{i+1}=2np_i$. Summing, $P(\overline{E_{i^*}})\le\sum_{4\le i<i^*}P(E_i\cap\overline{E_{i+1}})\le(i^*-4)/n^2$.

\emph{The last two levels.} Define $Y_t$ as before with $i=i^*$; now $np_{i^*}<6\ln n$, and Lemma~\ref{lem:chernoffdom} with $\mu=6\ln n$ gives $P(E_{i^*}\text{ and }\mu_{i^*+1}(n)\ge12\ln n)\le e^{-2\ln n}=n^{-2}$. Finally let $Z_t=1$ if ball $t$ has height at least $i^*+2$ and
$\nu_{i^*+1}(t-1)<12\ln n$. Given the past, $P(Z_t=1)\le(12\ln n/n)^2$, so $\mathbb E\sum_tZ_t\le144(\ln n)^2/n$, and by Markov's inequality $P(\sum_tZ_t\ge1)\le144(\ln n)^2/n$. If $\mu_{i^*+1}(n)<12\ln n$, then $\nu_{i^*+1}(t-1)<12\ln n$ for all $t$, and a ball of
height $i^*+2$ makes $\sum Z_t\ge1$. Altogether, the probability that some bin reaches load $i^*+2$ is at most
\[
\frac{i^*-4}{n^2}+\frac1{n^2}+\frac{144(\ln n)^2}n\le\frac{145(\ln n)^2}n,
\]
and otherwise the maximum load is at most $i^*+1\le\log_2\log_2n+6$.
\end{proof}

From $\ln n/\ln\ln n$ to $\log_2\log_2n$: an exponential improvement for the price of one more random choice. With $d$ choices the bound becomes $\ln\ln n/\ln d+O(1)$, so the big gain is from one to two; and Berthold V\"ocking showed in 1999 that splitting the bins into two halves, choosing one bin in each, and breaking ties
towards the left half, does better still. The principle, called \emph{the power of two choices}, is used in load balancers and in hash tables (cuckoo hashing).

\section{Hashing and data streams}\label{sec:streams}

A \emph{hash function} maps a large universe of keys into a small table of $w$ cells. No fixed function behaves well on every set of keys, so we choose the function at random from a family.

\begin{definition}
A family $\mathcal H$ of functions from a set $U$ to $\{0,\dots,w-1\}$ is \emph{universal} if $P_h\bigl(h(x)=h(y)\bigr)\le1/w$ for all $x\ne y$ in $U$, when $h$ is uniform in $\mathcal H$.
\end{definition}

\begin{proposition}[Carter and Wegman, 1979]\label{prop:universal}
Let $p$ be a prime, $U=\{0,\dots,p-1\}$ and $w\le p$. The functions $h_{a,b}(x)=\bigl((ax+b)\bmod p\bigr)\bmod w$, for $a\in\{1,\dots,p-1\}$ and $b\in\{0,\dots,p-1\}$, form a universal family.
\end{proposition}
\begin{proof}
Fix $x\ne y$. The map $(a,b)\mapsto(r,s)=(ax+b,ay+b)$ (in $\mathbb F_p$) is a bijection from the $p(p-1)$ pairs with $a\ne0$ onto the pairs $r\ne s$: given $r\ne s$, $a=(r-s)/(x-y)\ne0$ and $b=r-ax$. So $(r,s)$ is uniform among the pairs of distinct elements, and
$h(x)=h(y)$ means $r\equiv s\pmod w$. For each $r$, the number of $s\ne r$ in $\{0,\dots,p-1\}$ with $s\equiv r\pmod w$ is at most $\lceil p/w\rceil-1\le(p-1)/w$. So the probability of a collision is at most $\frac{p(p-1)/w}{p(p-1)}=\frac1w$.
\end{proof}

With a universal family, a hash table with chaining holding $n$ keys in $w$ cells answers each query in expected time $O(1+n/w)$, whatever the keys (Exercise~\ref{ex:hashchain}). Hashing is also the key to \emph{streaming}: data arrive as a sequence too long to store, the sequence of
packets through a router, of queries to a search engine, and we may keep only a small summary.

The first and smallest summary counts. Counting up to $n$ needs $\log_2n$ bits; in 1978 Robert Morris, at Bell Laboratories, needed many counters in very little memory and proposed to store only an approximation of the logarithm. The \emph{Morris counter} keeps an
integer $X$, initially $0$; at each event it increments $X$ with probability $2^{-X}$. The estimate of the number of events is $2^X-1$.

\begin{theorem}[Morris, 1978]\label{thm:morris}
After $n$ events, $\mathbb E[2^{X}]=n+1$ and $\operatorname{Var}(2^X)=n(n-1)/2$. Hence $P\bigl(|2^X-1-n|\ge\varepsilon n\bigr)\le1/(2\varepsilon^2)$, and $P\bigl(X\ge\log_2(n+1)+c\bigr)\le2^{-c}$: the counter uses $\log_2\log_2n+O(1)$ bits with high probability.
\end{theorem}
\begin{proof}
Let $X_n$ be the value after $n$ events. Given $X_n=x$, $\mathbb E[2^{X_{n+1}}]=2^x(1-2^{-x})+2^{x+1}2^{-x}=2^x+1$ and $\mathbb E[4^{X_{n+1}}]=4^x(1-2^{-x})+4^{x+1}2^{-x}=4^x+3\cdot2^x$. Taking expectations, $\mathbb E[2^{X_{n+1}}]=\mathbb E[2^{X_n}]+1$, so $\mathbb E[2^{X_n}]=n+1$, and
$\mathbb E[4^{X_{n+1}}]=\mathbb E[4^{X_n}]+3(n+1)$, so $\mathbb E[4^{X_n}]=1+\frac32n(n+1)$. Then $\operatorname{Var}(2^{X_n})=1+\frac32n(n+1)-(n+1)^2=\frac{n(n-1)}2$. The two inequalities are Chebyshev's (with $\operatorname{Var}\le n^2/2$) and Markov's, applied to $2^{X_n}$.
\end{proof}

The variance is large, but averaging $k$ independent counters divides it by $k$, and Exercise~\ref{ex:morrisbase} trades memory for accuracy differently. Philippe Flajolet's 1985 analysis of this counter, down to the fluctuations of its error, is a classic of the analysis of algorithms.

Next, frequencies. A stream of updates $(x,c)$, with $x$ in a universe $U$ and $c\ge0$, defines the frequencies $f_x$ (the total of the updates to $x$) and the total $N=\sum_xf_x$. The \emph{count--min sketch} of Graham Cormode and S.~Muthukrishnan (2005) keeps a $d\times w$ array of
counters, initially $0$, and $d$ hash functions $h_1,\dots,h_d$ chosen independently from a universal family; an update $(x,c)$ adds $c$ to the counters $(j,h_j(x))$, $1\le j\le d$, and the estimate of $f_x$ is $\hat f_x=\min_j\mathrm{CM}[j,h_j(x)]$.

\begin{theorem}[count--min]\label{thm:countmin}
For every $x$, $\hat f_x\ge f_x$, and $P(\hat f_x\ge f_x+\varepsilon N)\le\delta$ when $w=\lceil e/\varepsilon\rceil$ and $d=\lceil\ln(1/\delta)\rceil$. The sketch uses $O(\varepsilon^{-1}\log(1/\delta))$ counters, independently of $|U|$.
\end{theorem}
\begin{proof}
The counter $(j,h_j(x))$ holds $f_x+\Delta_j$ with $\Delta_j=\sum_{y\ne x}f_y\mathbf 1[h_j(y)=h_j(x)]\ge0$, so $\hat f_x\ge f_x$. By universality $\mathbb E\Delta_j\le N/w\le\varepsilon N/e$, and by Markov's inequality $P(\Delta_j\ge\varepsilon N)\le1/e$. The $d$ rows are independent, so all of
them overestimate by $\varepsilon N$ or more with probability at most $e^{-d}\le\delta$.
\end{proof}

Finally, a statistic of the whole stream: the \emph{second moment} $F_2=\sum_xf_x^2$, which measures how skewed the frequencies are (it is the self-join size of a database table). Here updates may have any sign. In 1996 Noga Alon, Yossi Matias and Mario Szegedy gave a sketch for it, the beginning of the
theory of streaming algorithms. Choose a random sign $s(x)\in\{\pm1\}$ for each key, from a \emph{$4$-wise independent} family: any four distinct keys get independent uniform signs. Such a family is easy to build: take a field $\mathbb F_q$ with $q=2^\ell\ge|U|$ elements (Tome~2, Theorem~9.6), a uniform random polynomial $g$ of degree at most $3$ over it,
and let $s(x)=(-1)^{b(g(x))}$ where $b$ is the first coordinate in a basis of $\mathbb F_q$ over $\mathbb F_2$. For distinct $x_1,\dots,x_4$ and any targets, exactly one polynomial of degree at most $3$ takes them (Lagrange interpolation), so $g(x_1),\dots,g(x_4)$ are independent and uniform, and so are their first coordinates.
The sketch maintains $Z=\sum_xs(x)f_x$: an update $(x,c)$ adds $s(x)c$. Storing $g$ takes four field elements.

\begin{theorem}[Alon, Matias and Szegedy, 1996]\label{thm:ams}
$\mathbb E[Z^2]=F_2$ and $\operatorname{Var}(Z^2)\le2F_2^2$. Consequently, the median of $m=\lceil8\ln(1/\delta)\rceil$ averages of $k=\lceil8/\varepsilon^2\rceil$ independent copies of $Z^2$ lies in $[(1-\varepsilon)F_2,(1+\varepsilon)F_2]$ with probability at least $1-\delta$.
\end{theorem}
\begin{proof}
Write $s_x=s(x)$. Then $Z^2=\sum_xf_x^2+\sum_{x\ne y}s_xs_yf_xf_y$, and $\mathbb E[s_xs_y]=\mathbb E s_x\,\mathbb Es_y=0$ for $x\ne y$, so $\mathbb E[Z^2]=F_2$. Expanding $Z^4=\sum s_xs_ys_zs_tf_xf_yf_zf_t$ over quadruples, a term has expectation $0$ as soon as some key appears an odd number of times (by
$4$-wise independence its sign is independent of the others and has mean $0$), and $1$ otherwise. The surviving quadruples are $(x,x,x,x)$, and those made of two distinct keys each appearing twice, in $\binom42=6$ arrangements. So
\[
\mathbb E[Z^4]=\sum_xf_x^4+6\sum_{x<y}f_x^2f_y^2,\qquad\operatorname{Var}(Z^2)=\mathbb E[Z^4]-F_2^2=4\sum_{x<y}f_x^2f_y^2\le2F_2^2.
\]
An average $A$ of $k$ independent copies has variance at most $2F_2^2/k$, and by Chebyshev $P(|A-F_2|\ge\varepsilon F_2)\le2/(k\varepsilon^2)\le\frac14$. The median of $m$ independent averages is bad only if at least half of them are; the number of good ones has mean at least
$\frac34m$, so by Hoeffding's inequality this has probability at most $e^{-2m/16}=e^{-m/8}\le\delta$.
\end{proof}

The space is $O(\varepsilon^{-2}\log(1/\delta))$ counters and random polynomials. Alon, Matias and Szegedy also proved that, by contrast, the maximum frequency $F_\infty$ requires space linear in $|U|$, by a reduction from communication complexity; their paper received the G\"odel prize in 2005. Estimating
the number of \emph{distinct} keys, $F_0$, in $O(\log\log|U|)$ bits per counter is the HyperLogLog algorithm of Chapter~11.

\section{Removing the coins}\label{sec:derandomise}

Some randomised algorithms can be made deterministic by replacing the coins by careful choices. A \emph{cut} $(S,V\setminus S)$ of a graph with $m$ edges, with $S$ chosen by fair independent coins, cuts each edge with probability $\frac12$, so its expected size is $m/2$: some cut has at least $m/2$ edges. This is the
probabilistic method of Tome~3, Chapter~12, and it hides an algorithm.

\begin{proposition}[the method of conditional expectations]\label{prop:condexp}
Process the vertices $v_1,\dots,v_n$ in order, putting $v_i$ on the side opposite to the majority of its already placed neighbours (either side in case of a tie). The resulting cut has at least $m/2$ edges.
\end{proposition}
\begin{proof}
Let $c_i$ be the expected size of the random cut where $v_1,\dots,v_i$ are placed as the algorithm placed them and the others by fair coins; $c_0=m/2$ and $c_n$ is the size of the final cut. Edges with an endpoint not yet placed contribute $\frac12$ each, whatever the placed vertices do. So placing $v_i$ changes the expectation only through
the edges from $v_i$ to placed neighbours, which contributed $\frac12$ each before and contribute $1$ or $0$ after: $c_i=c_{i-1}+\frac12(\#\text{cut}-\#\text{uncut})$ among these edges. The rule makes this change nonnegative, so $c_n\ge c_0=m/2$.
\end{proof}

The same reasoning applies whenever the conditional expectations can be computed, and pairwise independent families like the one of Proposition~\ref{prop:universal} reduce the number of random bits so much that all their values can be tried. How far derandomisation goes is one of the great open questions of complexity theory: Impagliazzo and Wigderson showed in 1997 that
if some problems solvable in exponential time need exponential-size circuits, then every polynomial randomised algorithm with bounded error can be made deterministic and polynomial ($\mathsf{BPP}=\mathsf P$). Most researchers believe that randomness, in the end, does not make polynomial time more powerful, but it certainly makes it simpler, and
the primality test of Chapter~10 was found decades before its deterministic successor.

\begin{beauty}
Each argument of this chapter is a few lines of probability, and each gives an algorithm that would be hard to find otherwise: a cut survives because most edges are elsewhere, a nonzero polynomial cannot hide from a random point, a second choice of bin collapses a logarithm to a double logarithm, and a sum of random signs
remembers a sum of squares. The coins do not replace thinking; they make it possible to think about the average over our own choices, where the adversary has no say.
\end{beauty}

\section{Exercises}

\begin{exercise}[$\star$]\label{ex:lvmc}
A Monte Carlo algorithm runs in time $T$ and is correct with probability at least $p$, and its answer can be checked in time $V$. Build a Las Vegas algorithm with expected time at most $(T+V)/p$.
\end{exercise}

\begin{exercise}[$\star\star$]\label{ex:cutsalpha}
Show that the cycle $C_n$ has exactly $\binom n2$ minimum cuts. Then show that for an integer $r\ge2$, a multigraph with minimum cut $k$ has at most $n^r$ cuts of size at most $rk/2$. (Contract down to $r$ vertices, then choose a random cut.)
\end{exercise}

\begin{exercise}[$\star$]\label{ex:szsharp}
Show that Lemma~\ref{lem:schwartzzippel} is sharp: for every $d\le|S|$ there is a polynomial of degree $d$ in one variable vanishing at exactly $d$ points of $S$. Show that in Proposition~\ref{prop:freivalds}, choosing $r$ uniform in $S^n$ for a finite set $S$ of the field lowers the error to $1/|S|$.
\end{exercise}

\begin{exercise}[$\star\star$]\label{ex:tuttefind}
Given a test that decides whether a graph has a perfect matching, find one with $m$ tests, where $m$ is the number of edges. With the test of Corollary~\ref{cor:lovasz}, how small must each error be for the whole procedure to succeed with probability $\frac12$?
\end{exercise}

\begin{exercise}[$\star$]\label{ex:stringeq}
Alice and Bob each hold a string of $n$ bits. Using Proposition~\ref{prop:fingerprint}, give a protocol in which Alice sends $O(\log n)$ bits to Bob and Bob decides whether the strings are equal, with error at most $1/n$. (Bertrand's postulate: there is a prime between $N$ and $2N$.)
\end{exercise}

\begin{exercise}[$\star$]\label{ex:emptybins}
With $n$ balls in $n$ bins, compute the expected number of empty bins and its limit divided by $n$. Show with the argument of Theorem~\ref{thm:maxload} that the number of empty bins is $n/e+O_P(\sqrt n)$.
\end{exercise}

\begin{exercise}[$\star$]\label{ex:hashchain}
Keys $x_1,\dots,x_n$ are stored in a table of $w$ cells by a hash function $h$ chosen from a universal family, each cell holding the list of its keys. Show that for every key $y$, the expected length of the list of $h(y)$ is at most $1+n/w$.
\end{exercise}

\begin{exercise}[$\star$]\label{ex:universalzero}
Show that if $a$ is allowed to be $0$ in Proposition~\ref{prop:universal}, the collision probability is still at most $1/w+1/p$, and that for $p=3$ and $w=2$ the family is no longer universal.
\end{exercise}

\begin{exercise}[$\star\star$]\label{ex:morrisbase}
Let $a>0$. A counter increments $X$ with probability $(1+a)^{-X}$. Show that $((1+a)^X-1)/a$ is an unbiased estimator of $n$ with variance $a\,n(n-1)/2$. What is the trade-off with memory?
\end{exercise}

\begin{exercise}[$\star\star$]\label{ex:maxcutlocal}
Show directly that the algorithm of Proposition~\ref{prop:condexp} cuts, for each $i$, at least half of the edges joining $v_i$ to earlier vertices, and deduce the bound $m/2$ again. On the path $abcd$ processed in the order $a,d,b,c$, with ties sent to the side of $a$, show that it does not find a maximum cut.
\end{exercise}

\section{Hints}

\begin{hint}{ex:lvmc}
Repeat run and check until the check passes. The number of rounds is geometric.
\end{hint}
\begin{hint}{ex:cutsalpha}
A minimum cut of $C_n$ removes two edges. For the second part, a fixed cut of size at most $rk/2$ survives the step from $i$ vertices with probability at least $1-r/i$.
\end{hint}
\begin{hint}{ex:szsharp}
$\prod_{s\in S_0}(x-s)$. For Freivalds, apply the lemma to the linear polynomial $(Dr)_i$ in $r$.
\end{hint}
\begin{hint}{ex:tuttefind}
Delete an edge if the graph still has a perfect matching without it.
\end{hint}
\begin{hint}{ex:stringeq}
Alice chooses a prime $p$ between $n^2$ and $2n^2$ and a random $x$, and sends $p$, $x$ and $F_x(a)$.
\end{hint}
\begin{hint}{ex:emptybins}
Linearity of expectation; for the variance, the indicators of empty bins are again negatively correlated.
\end{hint}
\begin{hint}{ex:hashchain}
Write the length as a sum of indicators and use linearity.
\end{hint}
\begin{hint}{ex:universalzero}
Condition on $a=0$ or $a\ne0$.
\end{hint}
\begin{hint}{ex:morrisbase}
Compute $\mathbb E[(1+a)^{X_{n+1}}\mid X_n]$ and $\mathbb E[(1+a)^{2X_{n+1}}\mid X_n]$ as in Theorem~\ref{thm:morris}.
\end{hint}
\begin{hint}{ex:maxcutlocal}
Count, at each step, the edges to placed neighbours; each edge is counted once, at its later endpoint.
\end{hint}

\section{Solutions}

\begin{solution}{ex:lvmc}
\idea A checkable Monte Carlo answer can be retried until it is certified.
\writeup Repeat: run the algorithm, check its answer, and stop when the check succeeds, outputting the answer. The output is always correct. Each round costs at most $T+V$ and succeeds with probability at least $p$ independently of the previous rounds, so the number of rounds $R$ satisfies
$P(R>r)\le(1-p)^r$ and $\mathbb ER=\sum_{r\ge0}P(R>r)\le1/p$. The expected time is at most $(T+V)/p$.
\end{solution}

\begin{solution}{ex:cutsalpha}
\idea Contraction stopped early keeps every small cut with probability polynomial in $1/n$.
\writeup In $C_n$ the minimum cut is $2$, and a cut of size $2$ is given by two distinct edges, whose removal leaves two arcs: these are the $\binom n2$ minimum cuts. For the second part, fix a cut $C$ of size at most $rk/2$. With $i>r$ vertices the multigraph has at least $ik/2$ edges, so $C$ survives
a contraction with probability at least $1-\frac{rk/2}{ik/2}=1-\frac ri$. Contracting down to $r$ vertices, $C$ survives with probability at least $\prod_{i=r+1}^n\frac{i-r}i=\binom nr^{-1}$. Then choose uniformly one of the $2^{r-1}-1$ cuts of the $r$-vertex multigraph; $C$ is among them, so the output is $C$ with probability at least
$\bigl(\binom nr2^{r-1}\bigr)^{-1}$. These events are disjoint for different $C$, so there are at most $\binom nr2^{r-1}\le\frac{n^r}{r!}2^{r-1}\le n^r$ such cuts, as $2^{r-1}\le r!$.
\end{solution}

\begin{solution}{ex:szsharp}
\idea Roots are the only obstruction.
\writeup For $S_0\subseteq S$ with $d$ elements, $f=\prod_{s\in S_0}(x-s)$ has degree $d$ and vanishes at exactly the $d$ points of $S_0$, so with probability $d/|S|$. In Freivalds' test with $D=AB-C$ and $d_{ij}\ne0$, $(Dr)_i=\sum_kd_{ik}r_k$ is a nonzero polynomial of degree $1$ in $r_1,\dots,r_n$, so by
Lemma~\ref{lem:schwartzzippel} it vanishes with probability at most $1/|S|$, and $Dr=0$ is no more likely.
\end{solution}

\begin{solution}{ex:tuttefind}
\idea Self-reducibility: a decision procedure finds the object by deleting what is not needed.
\writeup Go through the edges $e_1,\dots,e_m$ in turn, keeping a current graph $H$, initially $G$ (which has a perfect matching). If $H-e_i$ has a perfect matching, set $H=H-e_i$. At the end $H$ has a perfect matching, and every edge of $H$ is in all its perfect matchings (deleting it was refused, and later deletions only remove options); so $H$ is itself a
perfect matching, since a graph with a perfect matching $M$ and an extra edge $e\notin M$ would have allowed deleting $e$. With a test that errs with probability at most $\eta$ (the randomised test only errs by saying ``no''), the $m$ tests are all correct with probability at least $1-m\eta$ by the union bound, so $\eta\le1/(2m)$ suffices, for example
by taking $p\ge2nm$ in Corollary~\ref{cor:lovasz}, or by repeating each test $\lceil\log_2(2m)\rceil$ times.
\end{solution}

\begin{solution}{ex:stringeq}
\idea Send a short random evaluation, not the string.
\writeup Alice picks a prime $p$ with $n^2<p\le2n^2$ (one exists by Bertrand's postulate, and it can be found by trial division), a uniform $x\in\mathbb F_p$, and sends $p$, $x$ and $F_x(a)$, where the bits of $a$ are read as elements of $\mathbb F_p$: $O(\log n)$ bits. Bob computes $F_x(b)$ and answers ``equal'' if the fingerprints agree.
If $a=b$ he is right; if $a\ne b$ he errs with probability at most $(n-1)/p<1/n$ by Proposition~\ref{prop:fingerprint}.
\end{solution}

\begin{solution}{ex:emptybins}
\idea Each bin is empty with probability $(1-1/n)^n\approx1/e$; dependence between bins is negative.
\writeup The number of empty bins is $E=\sum_i\mathbf 1[L_i=0]$, so $\mathbb EE=n(1-\frac1n)^n$ and $\mathbb EE/n\to1/e$. Given $L_i=a$, $P(L_j=0)=(1-\frac1{n-1})^{n-a}$ increases with $a$; so $P(L_j=0\mid L_i=0)\le P(L_j=0\mid L_i>0)$, and as in Theorem~\ref{thm:maxload},
$P(L_i=0,L_j=0)\le P(L_i=0)P(L_j=0)$ (directly: $(1-\frac2n)^n\le(1-\frac1n)^{2n}$). Hence $\operatorname{Var}E\le\mathbb EE\le n$, and by Chebyshev $P(|E-\mathbb EE|\ge c\sqrt n)\le1/c^2$. Since $n(1-\frac1n)^n=n/e+O(1)$, the number of empty bins is $n/e+O_P(\sqrt n)$.
\end{solution}

\begin{solution}{ex:hashchain}
\idea Each other key lands in the same cell with probability at most $1/w$.
\writeup The list of $h(y)$ contains $y$ (if stored) and the keys $x_i\ne y$ with $h(x_i)=h(y)$. Its expected length is at most $1+\sum_{x_i\ne y}P(h(x_i)=h(y))\le1+n/w$ by universality. The bound holds for every set of keys: the probability is over $h$, not over the keys.
\end{solution}

\begin{solution}{ex:universalzero}
\idea The constant functions are rare, but they always collide.
\writeup Conditioning on $a$: $P(\text{collision})=P(a=0)\cdot1+P(a\ne0)\,P(\text{collision}\mid a\ne0)\le\frac1p+\frac1w$, by Proposition~\ref{prop:universal} for the second term. For $p=3$, $w=2$ and $x\ne y$: given $a\ne0$, the pair $(r,s)$ is uniform among the $6$ ordered pairs of distinct elements of $\{0,1,2\}$, and
$r\equiv s\pmod2$ only for $(0,2)$ and $(2,0)$, so the conditional probability is $\frac13$. Altogether $P(\text{collision})=\frac13+\frac23\cdot\frac13=\frac59>\frac12=\frac1w$.
\end{solution}

\begin{solution}{ex:morrisbase}
\idea Same computation with base $1+a$: a smaller base means more increments, a larger counter, and a smaller variance.
\writeup Given $X_n=x$, $\mathbb E[(1+a)^{X_{n+1}}]=(1+a)^x(1-(1+a)^{-x})+(1+a)^{x+1}(1+a)^{-x}=(1+a)^x+a$, so $\mathbb E[(1+a)^{X_n}]=1+an$ and the estimator is unbiased. Similarly $\mathbb E[(1+a)^{2X_{n+1}}\mid X_n=x]=(1+a)^{2x}+(2a+a^2)(1+a)^x$, so
$\mathbb E[(1+a)^{2X_n}]=1+(2a+a^2)\sum_{k<n}(1+ak)=1+(2a+a^2)\bigl(n+a\tfrac{n(n-1)}2\bigr)$, and $\operatorname{Var}((1+a)^{X_n})=a^3\frac{n(n-1)}2$ after simplification; dividing by $a^2$ gives the variance $a\frac{n(n-1)}2$. The counter reaches about $\log_{1+a}(an)\approx\frac1a\ln(an)$, so it needs about
$\log_2\log_2n+\log_2(1/a)$ bits: halving the variance costs one bit.
\end{solution}

\begin{solution}{ex:maxcutlocal}
\idea The rule is greedy and local: it guarantees half, not the optimum.
\writeup When $v_i$ is placed with $d_i$ already placed neighbours, it goes opposite to the majority of them (or anywhere in case of a tie), so at least $\lceil d_i/2\rceil$ of these $d_i$ edges are cut. Every edge is counted exactly once, at its later endpoint, so the cut has at least $\sum_id_i/2=m/2$ edges.
On the path $abcd$ in the order $a,d,b,c$: $a$ goes on side $1$; $d$ has no placed neighbour, a tie, so it goes on side $1$; $b$ has its neighbour $a$ on side $1$ and goes on side $2$; $c$ has neighbours $b$ (side $2$) and $d$ (side $1$), a tie, and goes on side $1$. The edges $ab$ and $bc$ are cut and $cd$ is not: $2$ edges, while
the alternating cut $\{a,c\},\{b,d\}$ cuts all $3$.
\end{solution}

\section*{Sources and further reading}

\begin{itemize}
\item \emph{Research papers to read.} D.~R.~Karger, \emph{Global min-cuts in RNC, and other ramifications of a simple min-cut algorithm}, Proceedings of the 4th ACM--SIAM SODA (1993), 21--30: read the contraction algorithm and its analysis first, which are Theorem~\ref{thm:karger}, then the counting of near-minimum cuts
(our Exercise~\ref{ex:cutsalpha}) and how Karger uses it for parallel algorithms. Y.~Azar, A.~Z.~Broder, A.~R.~Karlin and E.~Upfal, \emph{Balanced allocations}, SIAM Journal on Computing 29 (1999), 180--200 (STOC 1994): the proof of the upper bound is our Theorem~\ref{thm:twochoices} with sharper constants, and the paper also proves the matching lower bound and treats the
sequential and the dynamic versions.
\item \emph{The stories.} N.~Metropolis, \emph{The beginning of the Monte Carlo method}, Los Alamos Science 15 (1987), 125--130, for von Neumann's letter of March 1947, Ulam's solitaire and poker, and the naming; N.~Metropolis and S.~Ulam, \emph{The Monte Carlo method}, Journal of the American Statistical Association 44 (1949), 335--341;
M.~O.~Rabin, \emph{Probabilistic algorithms}, in \emph{Algorithms and Complexity: New Directions and Recent Results} (J.~F.~Traub, ed.), Academic Press, 1976, 21--39; R.~Solovay and V.~Strassen, \emph{A fast Monte-Carlo test for primality}, SIAM Journal on Computing 6 (1977), 84--85;
L.~Babai, \emph{Monte-Carlo algorithms in graph isomorphism testing}, Universit\'e de Montr\'eal, technical report 79-10, 1979, for the name Las Vegas.
\item D.~R.~Karger and C.~Stein, \emph{A new approach to the minimum cut problem}, Journal of the ACM 43 (1996), 601--640; R.~A.~DeMillo and R.~J.~Lipton, \emph{A probabilistic remark on algebraic program testing}, Information Processing Letters 7 (1978), 193--195; R.~Zippel, \emph{Probabilistic algorithms for sparse polynomials}, EUROSAM 1979,
Lecture Notes in Computer Science 72, 216--226; J.~T.~Schwartz, \emph{Fast probabilistic algorithms for verification of polynomial identities}, Journal of the ACM 27 (1980), 701--717; W.~T.~Tutte, \emph{The factorization of linear graphs}, Journal of the London Mathematical Society 22 (1947), 107--111; L.~Lov\'asz, \emph{On determinants, matchings, and random
algorithms}, Fundamentals of Computation Theory (FCT 1979), Akademie-Verlag, 565--574; K.~Mulmuley, U.~V.~Vazirani and V.~V.~Vazirani, \emph{Matching is as easy as matrix inversion}, Combinatorica 7 (1987), 105--113; R.~Freivalds, \emph{Probabilistic machines can use less running time}, IFIP Congress 1977, 839--842;
V.~Kabanets and R.~Impagliazzo, \emph{Derandomizing polynomial identity tests means proving circuit lower bounds}, Computational Complexity 13 (2004), 1--46.
\item B.~V\"ocking, \emph{How asymmetry helps load balancing}, Journal of the ACM 50 (2003), 568--589 (FOCS 1999); J.~L.~Carter and M.~N.~Wegman, \emph{Universal classes of hash functions}, Journal of Computer and System Sciences 18 (1979), 143--154; R.~Morris, \emph{Counting large numbers of events in small registers}, Communications of the ACM 21 (1978), 840--842;
P.~Flajolet, \emph{Approximate counting: a detailed analysis}, BIT 25 (1985), 113--134; G.~Cormode and S.~Muthukrishnan, \emph{An improved data stream summary: the count-min sketch and its applications}, Journal of Algorithms 55 (2005), 58--75; N.~Alon, Y.~Matias and M.~Szegedy, \emph{The space complexity of approximating the frequency moments},
Journal of Computer and System Sciences 58 (1999), 137--147 (STOC 1996); R.~Impagliazzo and A.~Wigderson, \emph{P = BPP if E requires exponential circuits: derandomizing the XOR lemma}, Proceedings of the 29th ACM STOC (1997), 220--229.
\item R.~Motwani and P.~Raghavan, \emph{Randomized Algorithms}, Cambridge University Press, 1995, and M.~Mitzenmacher and E.~Upfal, \emph{Probability and Computing}, 2nd edition, Cambridge University Press, 2017, are the textbooks; the second proves Theorem~\ref{thm:twochoices} in its Chapter~17.
\end{itemize}
